% Generated draft: requires mathematical and source review.
\documentclass[11pt]{article}
\usepackage[margin=25mm]{geometry}
\usepackage{fontspec}
\setmainfont{Latin Modern Roman}
\usepackage{amsmath,amssymb,mathtools,unicode-math}
\setmathfont{Latin Modern Math}
\usepackage{enumitem,needspace,xurl,hyperref}
\setlength{\parindent}{0pt}
\setlength{\parskip}{5pt}
\setlength{\emergencystretch}{3em}
\newcommand{\N}{\mathbb{N}}
\newcommand{\Z}{\mathbb{Z}}
\newcommand{\Q}{\mathbb{Q}}
\newcommand{\R}{\mathbb{R}}
\newcommand{\C}{\mathbb{C}}
\newcommand{\F}{\mathbb{F}}
\newcommand{\A}{\mathbb{A}}
\newcommand{\PP}{\mathbb{P}}
\newcommand{\E}{\mathbb{E}}
\newcommand{\eps}{\varepsilon}
\newcommand{\dd}{\,\mathrm{d}}
\begin{document}
\section*{UnsolvedMath ID: 10400205}
\textbf{Source problem code:} catalogue entry
\subsection{Definitions and mathematical statement}
Write \(Z_{\mathrm{KT}}\) for the Kuperberg--Thurston invariant of source reference [237], and \(Z_{\mathrm{LMO}}\) for the LMO invariant. The source asks two questions.
\begin{enumerate}
\item Find a surgery formula for \(Z_{\mathrm{KT}}\) in terms of the Chern--Simons series specified in Question 3.12. A surgery formula must compute \(Z_{\mathrm{KT}}(Y)\) from surgery-presentation data for \(Y\), using that particular series, and give an expression independent of changes of presentation representing the same admissible manifold. The series and its allowed inputs are not given here; no replacement series is specified.
\item Compare \(Z_{\mathrm{KT}}\) with \(Z_{\mathrm{LMO}}\) on their common intended domain, identifying the value spaces, normalization and auxiliary-data conventions needed for a precise relation. The source asks for a comparison and does not print an equality or other candidate formula.
\end{enumerate}
The manifold class, invariant definitions, surgery moves, coefficient or diagram spaces and completions are absent from the excerpt. Question 3.12 and the full entry for [237] are necessary to turn item (1) into a fixed formula target; they must not be inferred from the names alone.
\subsection{Short English statement}
Find a surgery expression for the Kuperberg--Thurston invariant using the specific Chern--Simons series of Question 3.12, and determine its precise relation to the LMO invariant.
\subsection{Sources}
\begin{itemize}
\item T. Ohtsuki, \emph{Problems on invariants of knots and 3-manifolds} (2002), Question 11.9, supplied PDF page 157. \url{https://msp.org/gtm/2002/04/gtm-2002-04-024s.pdf}. Recorded access 2026-07-29; extraction pdf-font-structure.
\item The source cites [237] for the Kuperberg--Thurston invariant and Question 3.12 for the Chern--Simons series. Neither the bibliography entry nor the referenced question is included in this record.
\item Local UnsolvedMath registry, \texttt{data/unsolvedmath/problems.json}, permanent ID \(10400205\), exact source code \texttt{catalogue entry}; research classification OPEN-TRIAGE.
\end{itemize}
\end{document}
